\section{Introduction}

Cyber-physical infrastructures increasingly require software-defined behavior at the edge, where constrained devices must execute application logic under strict reliability, security, timing, and life-cycle constraints. In industrial and research deployments alike, remote sites can include sensors, actuators, lightweight gateways, and mixed-criticality nodes that must be supervised continuously while remaining inexpensive, portable, and operationally manageable. Traditional embedded deployment strategies still rely heavily on manual flashing, board-specific provisioning, or tightly coupled update channels. Although workable at small scale, these approaches do not extend naturally to fleet-wide orchestration, and they provide limited observability to operators responsible for verifying which software version is actually running on which device at a given time.

In parallel, cloud-native ecosystems have matured around declarative orchestration, reconciliation loops, immutable artifacts, and service-level observability. Kubernetes has become the de facto orchestration substrate for distributed services because it separates desired state from observed state and continuously drives the latter toward the former through control loops and controllers~\cite{kubernetes}. This model has transformed data-center operations, but directly applying Kubernetes abstractions to constrained embedded endpoints remains difficult. Embedded devices often lack the resources to host container runtimes, the network reachability expected by cluster-native services, and the execution substrates assumed by modern cloud applications. Consequently, a direct transplant of cloud orchestration to the edge is usually infeasible.

WebAssembly (WASM) offers a promising execution abstraction for this gap. It enables portable, sandboxed workloads that can run across cloud, fog, and edge environments with a more stable artifact format than architecture-specific binaries~\cite{haas_wasm}. For embedded systems, WebAssembly Micro Runtime (WAMR) and Zephyr-compatible integrations provide a practical path toward microcontroller-oriented execution~\cite{wamr,zephyr}. However, runtime feasibility alone is not enough. The key unresolved problem is end-to-end orchestration: how cloud intent is expressed, how edge devices securely join the platform, how deployment commands reach constrained targets, and how execution state is reflected back into the control plane.

This paper introduces RETROSPECT, a full-stack platform that addresses this problem through a gateway-centric Cloud--Fog--Edge architecture. RETROSPECT combines Kubernetes CRDs, Rust-based control services, TLS-protected southbound communication, and CBOR-encoded protocol messages to manage device enrollment, secure connectivity, and WASM life-cycle operations across heterogeneous layers. The gateway mediates between the cloud control plane and constrained devices, allowing embedded nodes to remain lightweight while still participating in a declarative orchestration model.

The work is motivated by a concrete engineering observation: in cyber-physical environments, a secure enrollment exchange is necessary but insufficient. A practical platform must also maintain stable ownership of status updates, avoid false disconnection semantics during transient conditions, and expose enough evidence for operators to trust the control plane. This requirement turns the problem from a simple protocol design exercise into a system integration challenge spanning resource modeling, network communication, reconciliation logic, and deployment reproducibility.

The main contributions of the paper are as follows:
\begin{enumerate}
    \item A complete architecture and implementation of a Kubernetes-native edge orchestration platform for embedded WASM execution.
    \item A secure communication and enrollment workflow based on TLS and CBOR with gateway-mediated identity matching and explicit device association.
    \item A practical reconciliation strategy that links desired state and reported state for devices and applications while avoiding false disconnect loops.
    \item An end-to-end deployment and validation study on K3s with emulated embedded targets, including diagnosis and mitigation of a controller-level consistency defect.
    \item A reproducible artifact package that demonstrates how the full platform can be deployed, tested, and documented for peer-review-oriented evaluation.
\end{enumerate}

The remainder of the paper is structured as follows. Section~III discusses related work and positions the contribution. Section~IV formalizes the problem and design goals. Sections~V to VII present the architecture, protocol flow, and implementation details. Sections~VIII and IX describe the experimental methodology and results. Sections~X to XIII discuss implications, threats to validity, reproducibility, and future work.
